\begin{table}[htbp]
\centering
\small
\caption{Persistência espacial entre anos consecutivos.}
\label{tab:persistencia}
\begin{adjustbox}{max width=\linewidth}
\begin{tabular}{lrrrrr}
\toprule
& \multicolumn{2}{c}{Spearman} & & & \\
\cmidrule(lr){2-3}
Anos & Células 500 m & Incidência bairro & Jaccard top 10\% & Células top 10\% (\%) & Casos seguintes no top 10\% (\%) \\
\midrule
2013--2014 & 0,37 & 0,38 & 0,23 & 12,4 & 32,2 \\
2014--2015 & 0,38 & $-$0,17 & 0,11 & 12,2 & 18,8 \\
2015--2016 & 0,79 & 0,54 & 0,29 & 10,2 & 25,2 \\
2016--2017 & 0,45 & 0,26 & 0,27 & 10,2 & 26,5 \\
2017--2018 & 0,46 & 0,45 & 0,24 & 16,9 & 33,5 \\
2018--2019 & 0,58 & 0,44 & 0,20 & 13,7 & 28,0 \\
2019--2020 & 0,53 & 0,33 & 0,18 & 11,2 & 23,7 \\
2020--2021 & 0,51 & 0,38 & 0,24 & 12,9 & 26,3 \\
2021--2022 & 0,59 & 0,40 & 0,23 & 10,3 & 24,9 \\
2022--2023 & 0,46 & 0,38 & 0,30 & 12,4 & 30,8 \\
2023--2024 & 0,51 & 0,49 & 0,19 & 10,8 & 21,9 \\
2024--2025 & 0,60 & 0,30 & 0,20 & 11,3 & 23,4 \\
\midrule
Média & 0,52 & 0,35 & 0,22 & 12,0 & 26,2 \\
\bottomrule
\end{tabular}
\end{adjustbox}
\par\smallskip\begin{minipage}{0.95\linewidth}\footnotesize Células: contagem de casos por célula de 500 m. Incidência bairro: casos por habitante nos bairros com pelo menos 1.000 moradores. Top 10\%: células no decil superior de casos do ano anterior (empates incluídos).\end{minipage}
\end{table}
